\documentclass{ctexart}

% 常用宏包
\usepackage{amsmath,amssymb}
\usepackage{bm}
\usepackage{siunitx}
\usepackage{graphicx}
\usepackage{booktabs}
\usepackage{float}
\usepackage{caption}
\usepackage{subcaption}
\usepackage{geometry}
\usepackage{hyperref}

% 页面设置
\geometry{
    a4paper,
    left=2.5cm,
    right=2.5cm,
    top=2.5cm,
    bottom=2.5cm
}

\begin{document}

\begin{center}
    {\LARGE 《星系天文学》第3章作业}

    \vspace{1em}

    姓名：何飞杭 \qquad 学号：202631101066
\end{center}

\noindent\textbf{1. (1) 推导奥尔特公式，得到视向速度和切向速度随银经变化的关系以及奥尔特常数A、B。}

\vspace{15em}

\[V_r = V \cos \alpha - V_0 \sin l\]
其中
\[\cos \alpha = \frac{R_0}{R} \sin l\]
所以：
\[V_r = V \frac{R_0}{R} \sin l - V_0 \sin l = R_0(\frac{V}{R} - \frac{V_0}{R_0}) \sin l\]
------------------------------------------------------------------------------------------------------------------------------------
\[V_t = V \sin \alpha - V_0 \cos l\]
其中
\[\sin \alpha = \frac{R_0 \cos l-r}{R}\]
所以：
\[V_t = V \frac{R_0 \cos l - r}{R} - V_0 \cos l = R_0(\frac{V}{R} - \frac{V_0}{R_0}) \cos l - \frac{V}{R} r\]
------------------------------------------------------------------------------------------------------------------------------------
\[\frac{V}{R} = \frac{V_0}{R_0} + \frac{\mathrm{d}(V/R)}{\mathrm{d}R}(R-R_0) = \frac{V_0}{R_0} + \frac{1}{R_0^2}[R_0\frac{\mathrm{d}(V)}{\mathrm{d}R} - V_0](R-R_0)\]
对于 \(R \approx R_0 >> r\)
\[R = \sqrt{R_0^2 + r^2 - 2 R_0 r \cos l} \approx R_0 - r \cos l\]
所以：
\[\frac{V}{R} - \frac{V_0}{R_0} \approx \frac{1}{R_0^2}[R_0\frac{\mathrm{d}(V)}{\mathrm{d}R} - V_0](-r  \cos l) = -\frac{1}{R_0^2}[R_0\frac{\mathrm{d}(V)}{\mathrm{d}R} - V_0]r \cos l\]
------------------------------------------------------------------------------------------------------------------------------------
定义奥尔特常数：
\[
A=\frac{1}{2}\left(\frac{V_0}{R_0}-\frac{\mathrm{d}V}{\mathrm{d}R}\right)
\]
因此：
\[
\frac{1}{R_0^2}\left(R_0\frac{\mathrm{d}V}{\mathrm{d}R}-V_0\right)
=-\frac{2A}{R_0}
\]
代入前面的结果：
\[
\frac{V}{R}-\frac{V_0}{R_0}
\approx \frac{2A}{R_0}r\cos l
\]
代入视向速度得：
\[
V_r=2Ar\sin l\cos l
\]
利用二倍角公式：
\[
V_r=Ar\sin 2l
\]
------------------------------------------------------------------------------------------------------------------------------------
由于 \(R\approx R_0\)，因此：
\[
\frac{V}{R}\approx\frac{V_0}{R_0}
\]
代入切向速度得到：
\[
V_t=2Ar\cos^2 l-\frac{V_0}{R_0}r
\]
利用二倍角公式：
\[
V_t=Ar(1+\cos 2l)-\frac{V_0}{R_0}r
\]
定义第二个奥尔特常数：
\[
B=-\frac{1}{2}\left(\frac{V_0}{R_0}
+\frac{\mathrm{d}V}{\mathrm{d}R}\right)
\]
由 \(A\) 和 \(B\) 的定义可得：
\[
B=A-\frac{V_0}{R_0}
\]
因此：
\[
V_t=Ar\cos 2l+Br
\]
------------------------------------------------------------------------------------------------------------------------------------
综上，视向速度和切向速度随银经变化的关系为：
\[
V_r = Ar\sin 2l,\; V_t = Ar\cos 2l+Br
\]
其中，奥尔特常数 \(A\)、\(B\) 分别为：
\[
A=\frac{1}{2}\left(\frac{V_0}{R_0} -\frac{\mathrm{d}V}{\mathrm{d}R}\right),\; B=-\frac{1}{2}\left(\frac{V_0}{R_0} +\frac{\mathrm{d}V}{\mathrm{d}R}\right)
\]

\vspace{0.5em}

\noindent\textbf{(2) 根据太阳附近A、B的测量值判断该处是对应上升还是下降的旋转曲线？}

\vspace{0.25em}
\[m = -2.5 \log_{10}\frac{f_1 + f_2}{f_0} = -2.5 \log_{10}(10^{-0.4 m_1} + 10^{-0.4 m_2}) \approx 0.638\]
\vspace{0.5em}

\noindent\textbf{2. (1) 太阳(系)在轨道上绕着银河系中心公转一周的时间，称为银河年或宇宙年。请根据两个基本量R0、V0的数值计算银河年的大小。}

\vspace{0.25em}
\noindent 太阳到银河系中心的距离：
\[R_0 = 8.5\ \mathrm{kpc}\]
太阳附近的圆周旋转速度：
\[V_0 = 220\ \mathrm{km/s}\]
所以银河年为：
\[T = \frac{2 \pi R_0}{V_0} \approx 2.4 \times 10^8\ \mathrm{yr}\]
\vspace{0.5em}

\noindent\textbf{(2) 假设在太阳轨道半径处猎户臂密度波速是120km/s，求自太阳诞生以来46亿年太阳已穿行猎户臂多少次？}

\vspace{0.25em}
\noindent 相对速度为：
\[V_\mathrm{rel} = V_0 - V_\mathrm{arm} = 220 - 120 = 100\ \mathrm{km/s}\]
太阳穿行猎户臂的时间为：
\[T_\mathrm{cross} = \frac{2 \pi R_0}{V_\mathrm{rel}} \approx 5.22 \times 10^8\ \mathrm{yr}\]
所以太阳在46亿年中穿行猎户臂的次数为：
\[N = \frac{4.6 \times 10^9}{5.22 \times 10^8} \approx 8.81\]
因此太阳大约穿过猎户臂密度波8次
\vspace{0.5em}

\noindent\textbf{3. 距离为d Mpc的一个星系：\\
(1) 在地球上观测这个星系，证明星系 \(1''\) 区域对应于4.848d pc。}

\vspace{0.25em}
\noindent 当角度很小时：
\[\tan \theta \approx \theta\]
又因为：
\[\theta = 1'' = \frac{1}{3600} \times \frac{\pi}{180} \approx 4.848 \times 10^{-6}\ \mathrm{rad}\]
所以：
\[L = D \theta \approx 4.848d pc\]
\vspace{0.5em}

\noindent\textbf{(2) 利用视星等和流量、绝对星等和光度的关系，证明星系的表面亮度满足如下换算关系：
\[\boxed{
I(\mathrm{mag/arcsec^2})
=M_\odot+21.572-2.5\log_{10}L
}\]}

\vspace{0.25em}
\noindent 视星等与流量的关系为：
\[m = - 2.5 \log_{10} F + C\]
绝对星等与光度的关系为：
\[M - M_\odot = - 2.5 \log_{10} (\frac{L}{L_\odot})\]
因此，$1\,\mathrm{arcsec}^2$ 对应的实际面积为：
\[S = (4.848 d)^2\ \mathrm{pc}^2\]
设星系的面光度为 $L$，单位为
$L_\odot/\mathrm{pc}^2$，则该区域内的总光度为：
\[
L_{\mathrm{total}}=L(4.848d)^2
\]
根据绝对星等与光度的关系，该区域的绝对星等为：
\[
M=M_\odot-2.5\log_{10}
\left[L(4.848d)^2\right]
\]
距离模数为：
\[
m-M=5\log_{10}D-5
\]
其中 $D$ 以 $\mathrm{pc}$ 为单位，由于：
\[
D=d\times10^6\,\mathrm{pc}
\]
所以：
\[
m-M=5\log_{10}d+25
\]
代入得到：
\[
\begin{aligned}
m={}&M_\odot-2.5\log_{10}
\left[L(4.848d)^2\right]\\
&+5\log_{10}d+25
\end{aligned}
\]
展开对数：
\[
\begin{aligned}
m={}&M_\odot-2.5\log_{10}L\\
&-5\log_{10}(4.848)
-5\log_{10}d\\
&+5\log_{10}d+25
\end{aligned}
\]
化简可得：
\[
m=M_\odot+25-5\log_{10}(4.848)
-2.5\log_{10}L
\]
由于：
\[
25-5\log_{10}(4.848)=21.572
\]
又因为计算的是单位角面积内的视星等，
即星系的表面亮度 $I$，因此：
\[I(\mathrm{mag/arcsec^2})
=M_\odot+21.572-2.5\log_{10}L\]
\vspace{0.5em}

\noindent\textbf{(3) 如果测得该星系B波段面亮度\(I_B = 27 mag \cdot arcsec^-2\)，与太阳相比，该星系的一平方角秒发射了多少 B 波段的光？  （太阳相关数据：\(M_{⊙V} = 4.83，B-V=0.65\)）}

\vspace{0.25em}
\noindent 已知星系的 B 波段表面亮度为：
\[
I_B=27\ \mathrm{mag/arcsec^2}
\]
太阳的 V 波段绝对星等及颜色指数分别为：
\[
M_{\odot V}=4.83,\qquad B-V=0.65
\]
首先，根据颜色指数的定义：
\[
B-V=M_{\odot B}-M_{\odot V}
\]
得到太阳的 B 波段绝对星等：
\[
M_{\odot B}=M_{\odot V}+(B-V)
=4.83+0.65=5.48
\]
根据上一问得到的表面亮度与面光度之间的换算关系：
\[
I_B=M_{\odot B}+21.572
-2.5\log_{10}L_B
\]
将已知数据代入：
\[
27=5.48+21.572-2.5\log_{10}L_B
\]
整理得：
\[
2.5\log_{10}L_B=5.48+21.572-27
\]
即：
\[
\log_{10}L_B=\frac{0.052}{2.5}=0.0208
\]
因此：
\[
L_B=10^{0.0208}\approx1.05
\]
最终得到：
\[
L_B\approx1.05\ L_{\odot B}/\mathrm{pc}^2
\]
即该星系的 B 波段面光度约为
$1.05\ L_{\odot B}/\mathrm{pc}^2$，
表示每平方秒差距的实际面积所对应的 B 波段光度，
约为太阳 B 波段光度的 1.05 倍。

\end{document}